\section{Method}
This section details our methodology for long video generation. We begin by revisiting the conversion of bidirectional models into streaming autoregressive generators~\citep{yin2025causvid,huang2025self}. Building upon this, we introduce our novel strategies tailored for long-form video synthesis. The complete generative process is formalized in Algorithm~\ref{alg:alg1}.



\subsection{Background}
Video diffusion models, while powerful, typically require denoising along a multi-step noise schedule, which renders the generation process computationally intensive. A prevalent strategy to mitigate this computational burden is to distill the foundational model into a few-step generator. Prominent approaches in this domain include Distribution Matching (DM)~\citep{yin2024one,yin2024improved,luo2025learning} and Consistency Models (CM)~\citep{song2023consistency,wang2024phased}.
Building upon the methodologies of CausVid and Self-Forcing, we distill the original bidirectional teacher model into a few-step generator, then convert it into an autoregressive model. 
This conversion is accomplished by training a student model to replicate the Ordinary Differential Equation (ODE) trajectories sampled from the teacher. We refer to this procedure as an initialization stage (see ~\cref{sec.appendix.dmd} for implementation details).
The Self-Forcing method extends this approach by training the distilled model on self-generated rollouts of up to five seconds using techniques such as Distribution Matching Distillation (DMD) loss~\citep{yin2024improved}. While this technique effectively mitigates the over-exposure artifacts present in CausVid, it exhibits a critical limitation: a significant degradation in generative quality when producing sequences that exceed its constrained training horizon.


\subsection{Extend training beyond teacher's limit}
\label{sec.extended_dmd}

\paragraph{Motivation}
As discussed earlier, the teacher model is trained exclusively on five-second video segments. Consequently, distillation-based methods such as CausVid~\citep{yin2025causvid} and Self-Forcing~\citep{huang2025self} only enforce student-teacher distribution alignment within this limited temporal window. This constrained training objective leads to a precipitous decline in quality when generation extends beyond this five-second horizon.
Despite this performance collapse, we make a critical observation: \textbf{videos rolled out beyond the training horizon often retain structural coherence}, even if this coherence manifests as undesirable artifacts such as motion stagnation (a common failure mode in Self-Forcing). This suggests that the core problem is not a fundamental breakdown of the autoregressive mechanism, which correctly leverages the history KV cache to maintain context. Rather, the primary issue is the compounding of autoregressive errors during extended rollouts. These errors accumulate and eventually manifest as motion loss, scene freezing, and catastrophic degradation of visual fidelity. This insight motivates us to introduce a simple yet effective method to mitigate error accumulation, which is described in the following sections.

\paragraph{Backwards Noise Initialization}
A central challenge in extending student-teacher distillation to long-horizon video generation resides in the noise initialization strategy. In the short-horizon setting (i.e., for videos with a length up to $T$ frames), the student model can be directly supervised on complete trajectories sampled from the teacher, each originating from random noise. However, for long-horizon generation, a trajectory initialized from pure random noise is decoupled from the preceding video content, leading to a fundamental context misalignment since the sampled noise does not preserve the temporal dependencies of previously generated frames.
 Based on the observation mentioned above, we add noise back to the denoised latent vectors and use it as the starting noise which is also shown to boost the performance of distillation~\citep{yin2024improved}. While similar techniques of re-injecting noise have been employed in prior work~\citep{yin2024improved,yin2025causvid,huang2025self}, our motivation and application are distinct. Whereas they used this for short-video distillation, primarily to enhance single-shot quality or circumvent the need for real training data. We leverage it as a mechanism to enforce temporal consistency across long videos.
 Specifically, the student model is first rolled out to a sequence of $N$ clean frames, with $N \gg T$, where $T$ denotes the maximum horizon the teacher can reliably generate such as 5 seconds. We then re-inject noise into the student roll-out according to the same diffusion noise schedule $\{\sigma_t\}_{t=1}^N$. Formally, given the clean trajectory $\{x_t^{S}\}_{t=1}^N$ generated by the student, the generation is perturbed as: 
\begin{equation}
\label{eq:perturbation}
x_t = (1-\sigma_t)x_{0} + \sigma_t \epsilon,
\quad \text{where } 
x_{0} = x_{t-1} - \sigma_{t-1}\,\hat{\epsilon}_{\theta}(x_{t-1}, t-1) ,
\end{equation} where $\epsilon \sim \mathcal{N}(0, I)$ denotes Gaussian noise,  
and $\hat{\epsilon}_{\theta}$ is the noise prediction network parameterized by $\theta$.
which serves as the initial state for computing the teacher and student distributions. This approach ensures that the distribution divergence between student and teacher model is evaluated on trajectories that retain temporal consistency and correctly structured according to the prescribed noise schedule.

\begin{figure}[t]
\centering
\includegraphics[width=\linewidth]{main_workflow/ICLR_main_workflow.pdf}
\caption{Workflow between baselines and Self-Forcing++. Our method employ backward noise initialization, extended DMD and rolling KV Cache to effectively mitigates train-test discrepancies.}
\vspace{-5pt}
\label{fig:placeholder}
\end{figure}
\paragraph{Extended Distribution Matching Distillation} 
Our strategy for extending training to long videos is grounded in the observation that although the bidirectional teacher model is trained exclusively on short, five-second clips, it implicitly captures the underlying data distribution of the ``world'' from its training data. From this perspective, any short, contiguous video segment can be viewed as a sample from the marginal distribution of a valid, longer video sequence~\citep{bruce2024genie,alonso2024diffusion,li2025hunyuan}.
This intuition motivates our core methodological extension. Since  our baseline method Self-Forcing~\citep{ huang2025self} restricts the training duration to the first $T$ frames (typically $\sim$5 seconds), we instruct the student model to roll out to $N$ frames where $N \gg T$. We then uniformly sample a contiguous window of length $T$ from the generated sequence, and compute the distributional discrepancy between the student and teacher models within this window. This sliding-window distillation process is formalized as \cref{eq:dmd_extended}:
\begin{equation}
\footnotesize
\begin{aligned}
\nabla_{\theta}\mathcal{L}_{\substack{\text{DMD} \\ \text{extended}}}
&= \mathbb{E}_{t}\mathbb{E}_{z}\
   \Biggl[\nabla_{\theta}\operatorname{KL}\!\left(
         p^{S}_{\theta,t}(z)\,\Vert\,p^{T}_{t}(z)
      \right)
   \Biggr] \\
&\approx -\,\mathbb{E}_{t}\,\mathbb{E}_{\,i\sim \mathrm{Unif}\{1,\dots,N-K+1\}} \Biggl[
      \int\!\Bigl(
         s^{T}\!\bigl(\Phi(G_{\theta}(z_i),t),t\bigr)
         - s^{S}_{\theta}\!\bigl(\Phi(G_{\theta}(z_i),t),t\bigr)
      \Bigr)\,
      \frac{dG_{\theta}(z_i)}{d\theta}\,dz_i
   \Biggr],
\end{aligned}
\label{eq:dmd_extended}
\end{equation}

Here, $G_{\theta}(z)$ denotes the student generator rollout given a latent $z$, and $\Phi(\cdot,t)$ is the transformation process at timestep $t$.  
$p^{S}_{\theta,t}$ and $p^{T}_{t}$ represent the student and teacher distributions at time $t$, with corresponding scores $s^{S}_{\theta}$ and $s^{T}$.  
We uniformly sample a starting index $i \sim \mathrm{Unif}\{0,\dots,N-K\}$ from the student rollout of length $N$, and extract a window of length $K$.  
The student is then trained to minimize the average KL divergence between its distribution and the teacher’s distribution across this window.  
The window size $K$ is typically chosen to match the horizon to which the teacher model was originally trained to generate.

\textit{\textbf{Remark} Bi-directional diffusion can be seen as a process to gradually restore a degraded target in different denoising \textbf{time-steps}. Our method adapts the idea to autoregressive video generation regime by having a short-horizon teacher gradually restore student’s degraded rollouts at different temporal \textbf{time-frames} and then distills these correction knowledge back into the student model.}

\paragraph{Training with rolling KV Cache}

Despite using KV cache at inference time, CausVid~\citep{yin2025causvid}  still relies on recomputing overlapping frames and suffers from a severe over-exposure problem. Self-Forcing~\citep{huang2025self} attempts to address this but introduces a train-inference mismatch by using a fixed cache during training and a rolling cache at inference. Although this is partially mitigated by masking the first latent frame, the mismatch still leads to substantial error accumulation and temporal flickering in long videos (see~\cref{fig:main}).
In contrast, our method naturally eliminates this mismatch by employing a rolling KV cache during both training and inference. At training time, this cache is used to roll out sequences far beyond the teacher's supervisory horizon to compute the extended DMD as detailed above. Consequently, our approach greatly simplifies the entire process, requiring neither the recomputation of overlapping frames nor latent frame masking.

\subsection{Improving Long-Term Smoothness via GRPO}
\label{sec.grpo}
A common drawback of generative models ~\citep{xiao2023efficient,luo2025enhance} employing sliding-window or sparse attention mechanisms for long sequences generation is the gradual loss of long-term memory. This degradation often manifests as temporal inconsistencies such as objects abruptly emerging or vanishing or unnaturally rapid scene transitions. Although the method we proposed above has achieved strong results, we show that Group Relative Policy Optimization (GRPO), a reinforcement learning technique ~\citep{liu2025flow,xue2025dancegrpo}, can be utilized in autoregressive video generation framework when such phenomenon presents. The per step importance weight $\rho_{t,i} = \tfrac{\pi_{\theta}(a_{t,i}\,|\,s_{t,i})}{\pi_{\theta_{\mathrm{old}}}(a_{t,i}\,|\,s_{t,i})}$ where 
$\pi_{\theta}(a_{t,i}|s_{t,i})$ denotes the policy function for output $o_i$ at time step $t$ can be computed according to~\cref{eq:perturbation} and the overall generation probability can be computed as the sum of all the log probabilities in current autoregressive rollouts which we show in~\cref{appendix.sec.grpo}. 
To guide the optimization process towards temporally smooth outputs, we follow prior work~\citep{chefer2025videojam,nam2025optical} and use the relative magnitude of optical flow between consecutive frames as a proxy for motion continuity. 
\begin{algorithm}[H]
  \caption{Self-Forcing++ with Backward Noise Initialization (ours)}
  \small
  \begin{algorithmic}[1]
    \Require Student \(G_{\theta}\), teacher \(T_{\phi}\), cache size \(L\); 
             rollout length \(N \gg 5\)s; slice length \(K\) (5s);
             denoise steps \(\{t_1,\dots,t_T\}\)
    \Loop
      \State \(V \gets \mathrm{Rollout}(G_{\theta}, N, L)\)  
      \State Pick \(i \sim \{1,\dots,N{-}K{+}1\}\), set \(W \gets V[i:i{+}K{-}1]\) \Comment{uniform slice}
      \State Sample \(t \sim \{t_1,\dots,t_T\}\) \
      \State Backward noise initialization: 
             \(x_t(W) \gets \text{BackwardNoiseInit}(W, t)\) 
      \State \(\mathcal{L}_{\mathrm{DMD}} \gets 
             \mathrm{DMD}(G_{\theta}(x_t(W), t),~T_{\phi}(x_t(W), t))\)
      \State \(\theta \gets \theta - \eta \nabla_{\theta}\mathcal{L}_{\mathrm{DMD}}\)
    \EndLoop
    \State $R \gets \mathrm{OpticalFlowReward}(G_\theta)$;\quad $\theta \gets \mathrm{GRPO\_update}(\theta, R)$ 
  \end{algorithmic}
    \label{alg:alg1}
\end{algorithm}





\subsection{New metrics for long videos evaluation}
\label{sec.better_metrics}
Most prior works rely on VBench~\citep{huang2024vbench} to assess image and aesthetic quality in long video generation. We find, however, that outdated evaluation models make the benchmark favor over-exposed videos (e.g., CausVid) and degraded long videos (e.g., Self-Forcing), leading to inaccurate scores. To address this, we adopt Gemini-2.5-Pro~\citep{comanici2025gemini}, a state-of-the-art video MLLM with strong reasoning ability~\citep{chiang2024chatbot,liu2025videoreasonbench}. Our protocol defines key long-video issues such as over-exposure and error accumulation, prompts Gemini-2.5-Pro to rate videos along these axes, and aggregates the results onto a $0-100$ scale termed visual stability for consistent comparison. More details are provided in ~\cref{fig:vbench_problem} and~\cref{appendix.sec.gemini}.
\begin{figure}[h]
    \centering
    \includegraphics[width=\linewidth]{images/vbench_problem.pdf}
    \caption{The left figure shows the issue of image score issue and the right figure shows the issue of aesthetic score of regular and degraded images from earlier and later frames of the same video. VBench tends to overrate degraded and over-exposed frames rendering these two metrics unreliable.}
    \label{fig:vbench_problem}
    \vspace{-10pt}
\end{figure}